\documentclass[11pt]{article}

\usepackage[margin=1in]{geometry}
\usepackage[T1]{fontenc}
\usepackage[utf8]{inputenc}
\usepackage{lmodern}
\usepackage{amsmath,amssymb,mathtools}
\usepackage{booktabs,longtable,tabularx,array}
\usepackage{enumitem}
\usepackage[numbers,sort&compress]{natbib}
\usepackage[colorlinks=true,linkcolor=blue,citecolor=blue,urlcolor=blue]{hyperref}

\setlength{\parskip}{0.6em}
\setlength{\parindent}{0pt}
\setlist[itemize]{leftmargin=*}
\setlist[enumerate]{leftmargin=*}

\title{A Non-Expert Guide to the Stellarator Design Workflow}
\author{stellarator\_workflow repository draft}
\date{March 11, 2026}

\begin{document}
\maketitle

\begin{abstract}
This document explains, in plain language but with exact governing equations,
how a modern stellarator workflow turns a prescribed plasma boundary and profile
guesses into transport predictions, fusion power estimates, and design metrics.
The focus is the practical code chain built around VMEC or DESC for
equilibrium, BOOZ\_XFORM for Boozer coordinates, NEO / SFINCS / MONKES for
neoclassical and orbit-space metrics, GENE or GX for turbulence, and Trinity3D
or NEOPAX for profile evolution. The goal is not only to explain the physics
for non-specialists, but also to clarify the software interfaces that would be
needed to containerize the full workflow as a reproducible black box.
\end{abstract}

\tableofcontents

\section{Why this workflow exists}

Designing a stellarator is not a single computation. It is an iterative chain:

\begin{enumerate}
  \item choose a plasma boundary and rough pressure / current profile,
  \item solve the three-dimensional magnetohydrodynamic equilibrium,
  \item transform the equilibrium into coordinates suitable for transport and orbit calculations,
  \item compute neoclassical and turbulent transport,
  \item evolve the density and temperature profiles,
  \item feed the updated pressure and current back into the equilibrium calculation,
  \item evaluate whether the resulting device has acceptable confinement, transport, and fusion power.
\end{enumerate}

For a non-expert, the most important conceptual point is that the profile used
to \emph{start} the equilibrium solve is usually not the profile used to quote
the device's final performance. The final performance comes from the
transport-consistent profiles after the loop closes.

\section{What goes in, and what comes out}

We assume stellarator symmetry, so the boundary can be written using only
cosine $R_{mn}$ and sine $Z_{mn}$ terms:

\begin{equation}
R(\theta,\phi) = \sum_{m,n} R_{mn}\cos\!\left(m\theta - nN_{\mathrm{fp}}\phi\right),
\end{equation}

\begin{equation}
Z(\theta,\phi) = \sum_{m,n} Z_{mn}\sin\!\left(m\theta - nN_{\mathrm{fp}}\phi\right).
\end{equation}

In VMEC input files these are usually represented as \texttt{RBC(m,n)} and
\texttt{ZBS(m,n)}. In DESC the same physical surface is stored in its own
spectral basis.

The remaining inputs are profile guesses:

\begin{itemize}
  \item a pressure profile $p$,
  \item either rotational transform $\iota$ or a current profile,
  \item a total toroidal flux / magnetic-field scale.
\end{itemize}

The converged output of the full workflow is not just ``an equilibrium''. It is
the combination of:

\begin{itemize}
  \item geometry metrics: major radius $R$, minor radius $a$, aspect ratio $R/a$, volume, $\beta$, $\iota$,
  \item Boozer-space metrics: symmetry-breaking harmonics, $I(\psi)$, $G(\psi)$, trapped-particle proxies,
  \item neoclassical metrics: $\epsilon_{\mathrm{eff}}$, bootstrap current, ambipolar $E_r$, neoclassical particle and heat fluxes,
  \item turbulence metrics: linear growth rates, nonlinear heat fluxes, stiffness,
  \item whole-device metrics: transport-consistent $n(r)$ and $T(r)$, fusion power, auxiliary power, and $Q=P_{\mathrm{fus}}/P_{\mathrm{aux}}$.
\end{itemize}

When the final species profiles are available, the fusion power is computed from
the transport-consistent pressure and composition:

\begin{equation}
P_{\mathrm{fus}} = \int dV \; n_D n_T \langle \sigma v \rangle_{DT} E_{DT},
\end{equation}

with $E_{DT}\approx 17.6$ MeV and the thermal reactivity
$\langle \sigma v \rangle_{DT}$ typically evaluated using the Bosch-Hale
parametrization \citep{BoschHale1992}.

\section{The workflow as a data pipeline}

Table~\ref{tab:contracts} summarizes the software contracts in the order they
usually appear in a design loop.

\begin{longtable}{p{0.13\linewidth}p{0.19\linewidth}p{0.21\linewidth}p{0.21\linewidth}p{0.19\linewidth}}
\caption{Core data contracts in a stellarator design workflow.}
\label{tab:contracts}\\
\toprule
Stage & Codes & Typical inputs & Typical outputs & Downstream role / direct objective \\
\midrule
\endfirsthead
\toprule
Stage & Codes & Typical inputs & Typical outputs & Downstream role / direct objective \\
\midrule
\endhead
Equilibrium &
VMEC, vmec\_jax, DESC &
boundary Fourier coefficients, pressure coefficients, iota or current coefficients, toroidal flux &
$wout$ or HDF5 equilibrium, $R$, $a$, aspect ratio, volume, $\beta$, $\iota$, field harmonics &
feeds everything downstream; geometry metrics are also direct objectives \\

Boozer transform &
BOOZ\_XFORM, booz\_xform\_jax, DESC Boozer utilities &
VMEC- or DESC-like equilibrium &
\texttt{boozmn}, Boozer $B_{mn}$, Boozer geometry, $I(\psi)$, $G(\psi)$, angle shifts &
feeds NEO, SFINCS, MONKES, some gyrokinetic geometry modules; symmetry-breaking spectrum is a direct objective \\

Orbit / ripple diagnostics &
NEO, NEO\_JAX &
Boozer Fourier data, surface list, control file &
$\epsilon_{\mathrm{eff}}^{3/2}$, trapped-particle class contributions, reference radii &
mostly a screening metric or optimization objective; usually not directly advanced by a transport solver \\

Neoclassical transport &
SFINCS, sfincs\_jax &
Boozer geometry, species gradients, $E_r$ guess, collision model &
particle flux, heat flux, bootstrap current, $E_r$, $\Phi_1$, transport matrices &
can feed a transport solver directly; bootstrap current and pressure updates can be sent back to equilibrium \\

Reduced neoclassical stage &
MONKES, NEOPAX &
Boozer or DESC geometry, Maxwellians, $E_r$, monoenergetic scan database &
$D_{ij}(r,\nu,E_r)$, thermal $L_{ij}$ matrices, ambipolar $E_r$, reduced transport evolution &
MONKES feeds NEOPAX; both can be used in optimization loops when full SFINCS is too expensive \\

Turbulence &
GENE, GENE-3D, GX, SPECTRAX-GK &
geometry, species profiles and gradients, collisions, electromagnetic parameters &
growth rates, frequencies, nonlinear heat and particle fluxes, spectra &
heat fluxes are both direct objectives and inputs to global transport solvers \\

Global transport &
Trinity3D, NEOPAX &
initial profiles, source models, geometry, neoclassical and turbulent fluxes &
updated $n(r)$, $T(r)$, $p(r)$, $E_r(r)$, fusion-power metrics, profile histories &
closes the loop by producing the new pressure and often new current / bootstrap information \\
\bottomrule
\end{longtable}

Two practical distinctions are worth making explicitly:

\begin{itemize}
  \item $\epsilon_{\mathrm{eff}}$ is extremely useful for optimization, but it is not normally the quantity that Trinity3D advances.
  \item nonlinear heat flux from GX or GENE is both a design objective and a direct transport input.
\end{itemize}

\section{Code-by-code details}

\subsection{VMEC and vmec\_jax: ideal-MHD equilibrium}

\paragraph{What these codes do.}
VMEC solves the three-dimensional ideal-MHD equilibrium for a prescribed
boundary and prescribed pressure / current data. The JAX implementation
\texttt{vmec\_jax} is designed to preserve the same physics and very similar
input / output contracts while enabling differentiation and JIT execution
\citep{VMECRepo,VMECJaxRepo}.

\paragraph{Governing equations.}
The equilibrium is defined by
\begin{equation}
\nabla p = \mathbf{J}\times\mathbf{B},
\qquad
\nabla\cdot\mathbf{B}=0,
\qquad
\mathbf{J}=\frac{1}{\mu_0}\nabla\times\mathbf{B}.
\end{equation}

VMEC uses straight-field-line coordinates. Introducing
\begin{equation}
u = \theta + \lambda(s,\theta,\zeta),
\qquad
\frac{du}{d\zeta} = \iota(s),
\end{equation}
the equilibrium is found as a stationary point of the ideal-MHD energy
functional \citep{HirshmanWhitson1983}:
\begin{equation}
W = \frac{1}{(2\pi)^2}\int \left(\frac{B^2}{2} + \frac{p}{\gamma-1}\right)dV.
\end{equation}

In VMEC's flux-coordinate representation, the contravariant field components
used in the implementation are \citep{VMECJaxRepo}:
\begin{equation}
B^\zeta = \frac{\Phi'(s) + \mathrm{lamscale}\,\partial_\theta \lambda}
{\mathrm{signgs}\,\sqrt{g}\,2\pi},
\end{equation}
\begin{equation}
B^\theta = \frac{\chi'(s) - \mathrm{lamscale}\,\partial_\zeta \lambda}
{\mathrm{signgs}\,\sqrt{g}\,2\pi}.
\end{equation}

\paragraph{Inputs.}
The legacy VMEC input contract is typically:
\begin{itemize}
  \item boundary coefficients \texttt{RBC(m,n)} and \texttt{ZBS(m,n)},
  \item pressure coefficients \texttt{AM},
  \item either iota coefficients \texttt{AI} or current coefficients \texttt{AC},
  \item a flux scale such as \texttt{PHIEDGE},
  \item grid / resolution / convergence controls.
\end{itemize}

\paragraph{Outputs.}
The standard product is \texttt{wout\_*.nc}, which contains the geometry and
diagnostic scalars needed by downstream codes: aspect ratio, major radius,
minor radius, volume, total beta, iota, pressure profiles, field harmonics,
Jacobian information, and real-space geometry on flux surfaces.

\paragraph{How these outputs are used.}
The equilibrium is the universal upstream state. BOOZ\_XFORM uses it directly,
GX geometry utilities derive flux-tube metrics from it, and transport solvers
use its flux-surface geometry. Many device objectives, such as target aspect
ratio or target volume, are already available at this stage.

\subsection{DESC: differentiable pseudo-spectral equilibrium}

\paragraph{What DESC does.}
DESC solves the same physical equilibrium problem, but in an inverse
pseudo-spectral formulation designed for optimization and automatic
differentiation \citep{DudtKolemen2020,Panici2023,Conlin2023,Dudt2023QS,DESCRepo}.

\paragraph{Governing equations.}
DESC represents the magnetic field in flux coordinates $(\rho,\theta,\zeta)$ as
\begin{equation}
\mathbf{B}
=
\frac{\partial_\rho \psi}{2\pi\sqrt{g}}
\left[
\left(\iota - \frac{\partial \lambda}{\partial \zeta}\right)\mathbf{e}_\theta
+
\left(1 + \frac{\partial \lambda}{\partial \theta}\right)\mathbf{e}_\zeta
\right].
\end{equation}

The current density follows from Ampere's law:
\begin{equation}
J^\rho = \frac{\partial_\theta B_\zeta - \partial_\zeta B_\theta}{\mu_0\sqrt{g}},
\qquad
J^\theta = \frac{\partial_\zeta B_\rho - \partial_\rho B_\zeta}{\mu_0\sqrt{g}},
\qquad
J^\zeta = \frac{\partial_\rho B_\theta - \partial_\theta B_\rho}{\mu_0\sqrt{g}}.
\end{equation}

Force balance is
\begin{equation}
\mathbf{F}\equiv\mathbf{J}\times\mathbf{B}-\nabla p = 0.
\end{equation}
In practice DESC solves a collocated nonlinear system in radial and helical
components of the force residual.

\paragraph{Inputs and outputs.}
DESC accepts its own input format or imported VMEC inputs. It requires the
boundary, pressure coefficients, iota or current coefficients, and toroidal
flux. It outputs an HDF5 equilibrium together with a wide set of geometry and
optimization diagnostics.

\paragraph{How DESC fits in the workflow.}
DESC can replace VMEC as the equilibrium engine, and it is especially useful
when the workflow is part of a differentiable optimization loop. DESC can also
compute several geometry objectives directly, reducing the number of external
post-processing steps needed during optimization.

\subsection{BOOZ\_XFORM and booz\_xform\_jax: VMEC to Boozer coordinates}

\paragraph{What these codes do.}
Boozer coordinates are the standard coordinate system for much of stellarator
neoclassical analysis. BOOZ\_XFORM converts a VMEC-like equilibrium into the
Boozer representation \citep{Boozer1983,BoozXformRepo,BoozXformJaxRepo}.

\paragraph{Governing equations.}
Boozer coordinates satisfy
\begin{equation}
\mathbf{B}
=
\nabla\psi\times\nabla\theta_B
+
\iota(\psi)\nabla\zeta_B\times\nabla\psi
=
\beta\nabla\psi + I(\psi)\nabla\theta_B + G(\psi)\nabla\zeta_B.
\end{equation}

Let $\zeta_0$ and $\theta_0$ denote the original angle coordinates and define
\begin{equation}
\zeta_B = \zeta_0 + \nu.
\end{equation}
Then the Boozer poloidal angle is
\begin{equation}
\theta_B = \theta_0 + \lambda + \iota\nu.
\end{equation}

The angle-shift field $\nu$ is determined by the covariant field components:
\begin{equation}
B_{\theta_0}
=
\left(1+\frac{\partial\lambda}{\partial\theta_0}\right)I
+
(G+\iota I)\frac{\partial\nu}{\partial\theta_0},
\end{equation}
\begin{equation}
B_{\zeta_0}
=
G + I\frac{\partial\lambda}{\partial\zeta_0}
+
(G+\iota I)\frac{\partial\nu}{\partial\zeta_0}.
\end{equation}

The implementation summarized in the BOOZ\_XFORM theory note expresses this as
\begin{equation}
\nu = \frac{w - I\lambda}{G+\iota I},
\end{equation}
where the Fourier coefficients of $w$ are reconstructed from the original
covariant field harmonics.

\paragraph{Inputs and outputs.}
The standard inputs are a VMEC \texttt{wout} file and a requested Boozer
resolution. The standard output is \texttt{boozmn\_*.nc}, containing Boozer
harmonics such as:
\begin{itemize}
  \item \texttt{bmnc\_b}: $|B|$ in Boozer coordinates,
  \item \texttt{rmnc\_b}, \texttt{zmns\_b}: geometry in Boozer coordinates,
  \item \texttt{pmns\_b}: toroidal angle shift,
  \item \texttt{gmn\_b}: Jacobian harmonics,
  \item \texttt{buco\_b}, \texttt{bvco\_b}: Boozer currents.
\end{itemize}

\paragraph{How these outputs are used.}
NEO and SFINCS typically expect Boozer-space data. The Boozer spectrum is also
one of the cleanest direct measures of quasi-symmetry.

\subsection{NEO and NEO\_JAX: effective ripple in the $1/\nu$ regime}

\paragraph{What these codes do.}
NEO evaluates effective helical ripple and related trapped-particle quantities
from Boozer-coordinate equilibria \citep{Nemov1999,NEORepo,NEOJaxRepo}.

\paragraph{Governing equations.}
The field-line ODE accumulates the integrals
\begin{equation}
y_2 = \int d\phi\,B^{-2},
\qquad
y_3 = \int d\phi\,|\nabla\psi|B^{-2},
\qquad
y_4 = \int d\phi\,K_G B^{-3},
\end{equation}
along with pitch-angle dependent trapped-particle integrals such as
\begin{equation}
I_f = \int d\phi\,\sqrt{1-\frac{B}{B_0\eta}}\,B^{-2},
\end{equation}
\begin{equation}
H_f = \int d\phi\,\sqrt{1-\frac{B}{B_0\eta}}
\left(\frac{4}{B/B_0} - \frac{1}{\eta}\right)\frac{K_G}{\sqrt{\eta}}B^{-2}.
\end{equation}

For each trapped-particle class,
\begin{equation}
\epsilon_{\mathrm{eff}}^{3/2}(m)
=
C_\epsilon \frac{y_2}{y_3^2}\,\mathrm{BigInt}(m),
\qquad
C_\epsilon=\frac{\pi R_0^2 \Delta\eta}{8\sqrt{2}}.
\end{equation}
The total \texttt{epstot} is the sum over classes.

\paragraph{Inputs and outputs.}
The usual inputs are \texttt{boozmn\_*.nc} and a \texttt{neo\_in.*} control
file. The output \texttt{neo\_out.*} contains $\epsilon_{\mathrm{eff}}$,
surface radii, iota, and optional current-related diagnostics.

\paragraph{How these outputs are used.}
These quantities are often used as optimization metrics or rapid screening
diagnostics. They usually do not directly enter a profile evolution equation in
Trinity3D.

\subsection{SFINCS and sfincs\_jax: full neoclassical drift-kinetic transport}

\paragraph{What these codes do.}
SFINCS solves the stellarator neoclassical drift-kinetic equation with realistic
collisions and, optionally, a poloidally varying electrostatic potential
\citep{Landreman2014SFINCS,Mollen2018,SFINCSRepo,SFINCSJaxRepo}.

\paragraph{Governing equations.}
One common form of the first-order kinetic equation solved by SFINCS is
\begin{align}
\left(
v_\parallel\mathbf{b}
+
\frac{d\Phi_0}{dr}\frac{\mathbf{B}\times\nabla r}{B^2}
\right)\cdot\nabla f_{s1}
+
\left[
-\frac{(1-\xi^2)v}{2B}\nabla_\parallel B + \cdots
\right]\frac{\partial f_{s1}}{\partial \xi}
- (\mathbf{v}_{ms}\cdot\nabla r)\frac{Z_s e}{2T_s x_s}\frac{d\Phi_0}{dr}\frac{\partial f_{s1}}{\partial x_s}
\nonumber\\
+
(\mathbf{v}_{ms}\cdot\nabla r)
\left[
\frac{1}{n_s}\frac{dn_s}{dr}
+
\frac{Z_s e}{T_s}\frac{d\Phi_0}{dr}
+
\left(x_s^2-\frac{3}{2}\right)\frac{1}{T_s}\frac{dT_s}{dr}
\right]f_{sM}
=
C_s[f_{s1}] + S_s.
\end{align}

The linearized collision operator is
\begin{equation}
C_s[f_s] = \sum_b C_{sb}^l[f_s,f_b],
\qquad
C_{ab}^l = C_{ab}^L + C_{ab}^E + C_{ab}^F.
\end{equation}

When $\Phi_1$ is included, the system is coupled to quasineutrality:
\begin{equation}
\lambda + \sum_s Z_s\int d^3v\,(f_{s0}+f_{s1}) = 0,
\end{equation}
together with moment and gauge constraints such as $\langle\Phi_1\rangle=0$.

\paragraph{Inputs.}
SFINCS typically needs:
\begin{itemize}
  \item Boozer geometry (from \texttt{boozmn} or \texttt{.bc}),
  \item species masses, charges, densities, temperatures, and radial gradients,
  \item a radial electric-field guess,
  \item collision-operator and numerical-resolution choices.
\end{itemize}

\paragraph{Outputs.}
The standard output file \texttt{sfincsOutput.h5} includes particle flux, heat
flux, bootstrap current, flow moments, transport matrices, and optional
$\Phi_1$ diagnostics.

\paragraph{How these outputs are used.}
In a reduced workflow, SFINCS can directly provide the neoclassical fluxes used
in a global transport solve. Bootstrap current and ambipolar $E_r$ can also be
fed back to the equilibrium stage.

\subsection{MONKES: monoenergetic neoclassical coefficients}

\paragraph{What MONKES does.}
MONKES is a fast solver for monoenergetic neoclassical transport coefficients,
designed to play a role similar to DKES but with improved modern numerics
\citep{Escoto2024MONKES,MONKESRepo}.

\paragraph{Governing equations.}
In the JAX port used here, the Legendre-expanded pitch-angle operators are
\begin{equation}
L_k(f)=\frac{k}{2k-1}\left[\mathbf{b}\cdot\nabla f + \frac{k-1}{2}\frac{\mathbf{b}\cdot\nabla B}{B}f\right],
\end{equation}
\begin{equation}
U_k(f)=\frac{k+1}{2k+3}\left[\mathbf{b}\cdot\nabla f - \frac{k+2}{2}\frac{\mathbf{b}\cdot\nabla B}{B}f\right],
\end{equation}
\begin{equation}
D_k(f)= -\frac{\hat E_r}{\psi_r\langle B^2\rangle}\mathbf{B}\times\nabla\psi\cdot\nabla f + \frac{k(k+1)}{2}\hat\nu f.
\end{equation}

The monoenergetic coefficients are then assembled from flux-surface averages of
the solved response:
\begin{equation}
D_{ij} =
\begin{pmatrix}
D_{11} & D_{12} & D_{13}\\
D_{21} & D_{22} & D_{23}\\
D_{31} & D_{32} & D_{33}
\end{pmatrix}.
\end{equation}

\paragraph{Inputs and outputs.}
The code takes a single-surface field, species Maxwellians, a radial electric
field, a speed, and numerical resolution. The main output is $D_{ij}$, together
with the perturbed distribution and source vectors.

\paragraph{How MONKES fits in the workflow.}
Its main downstream role is the creation of the database
$D_{ij}(r,\nu,E_r)$ for NEOPAX.

\subsection{GENE, GENE-3D, GX, and SPECTRAX-GK: turbulent transport}

\paragraph{What these codes do.}
These codes solve the gyrokinetic equations, which are the standard reduced
kinetic model for low-frequency plasma turbulence. GENE and GENE-3D are
high-fidelity grid-based Eulerian tools; GX uses a GPU-native
Fourier--Hermite--Laguerre representation \citep{Gorler2011GENE,Merz2012GENE,Xanthopoulos2020GENE3D,Mandell2018,Mandell2024GX,GENEWebsite,GXRepo,SpectraxGKRepo}.

\paragraph{Governing equations.}
A generic $\delta f$ gyrokinetic equation can be written
\begin{equation}
\frac{\partial h_s}{\partial t}
+
v_\parallel\mathbf{b}\cdot\nabla h_s
+
\mathbf{v}_{Ds}\cdot\nabla h_s
+
\mathbf{v}_E\cdot\nabla h_s
-
C[h_s]
=
-\frac{Z_s e F_{Ms}}{T_s}\frac{\partial\langle\chi\rangle}{\partial t}
-
\mathbf{v}_\chi\cdot\nabla F_{Ms},
\end{equation}
closed by quasineutrality and, for electromagnetic calculations, the
appropriate field equations.

GX represents velocity space spectrally:
\begin{equation}
h_s =
\sum_{\ell,m,k_x,k_y}\hat h_{s,\ell,m}(z,t)
e^{i(k_x x + k_y y)}
H_m\!\left(\frac{v_\parallel}{v_{ts}}\right)
L_\ell\!\left(\frac{v_\perp^2}{v_{ts}^2}\right)F_{Ms}.
\end{equation}

\paragraph{Inputs.}
The normal inputs are:
\begin{itemize}
  \item magnetic geometry, commonly derived from VMEC,
  \item local profiles and gradients,
  \item collisionality, beta, flow shear, and numerical resolution.
\end{itemize}

\paragraph{Outputs.}
These codes provide growth rates, mode frequencies, nonlinear heat and particle
fluxes, fluctuation spectra, and field energies. For workflow purposes the most
important outputs are the saturated heat and particle fluxes.

\paragraph{How these outputs are used.}
Those fluxes are precisely what a transport solver needs. They are also
frequently used directly as optimization objectives. This is why GX is so
useful inside optimization loops.

\subsection{Trinity3D: whole-device profile evolution}

\paragraph{What Trinity3D does.}
Trinity3D advances radial density and temperature profiles using fluxes supplied
by turbulence and neoclassical submodels \citep{TrinityDocs}.

\paragraph{Governing equations.}
In continuous form, the model is a 1D system of conservation laws:
\begin{equation}
\frac{\partial n_s}{\partial \tau}
+
\mathcal{G}(\rho)\frac{\partial F_{n,s}}{\partial \rho}
=
S_{n,s},
\end{equation}
\begin{equation}
\frac{\partial p_s}{\partial \tau}
+
\mathcal{G}(\rho)\frac{\partial F_{p,s}}{\partial \rho}
=
\frac{2}{3}S_{p,s}.
\end{equation}

The current implementation advances a stacked profile vector $y$ with an
implicit linearized update:
\begin{equation}
\left(d_1 I + \alpha\Psi\right)y^{m+1}
=
-d_0 y^m - d_{-1}y^{m-1}
+
\alpha\Psi y^m
-
\alpha\left[G(F^+ - F^-) - S\right]
-
(1-\alpha)\left[G(F^+_m - F^-_m) - S_m\right].
\end{equation}

\paragraph{Inputs and outputs.}
The code reads a TOML input specifying geometry, species profiles, source
models, and transport submodels. In stellarator examples, the geometry is often
VMEC, turbulence comes from GX, and neoclassical fluxes can come from SFINCS.
The outputs are updated profiles, flux histories, and the coupled sub-run
artifacts.

\paragraph{How Trinity3D fits in the workflow.}
This is typically the stage that turns local transport coefficients into a
global profile prediction and hence into an updated pressure profile for the
next equilibrium iteration.

\subsection{NEOPAX: a reduced JAX neoclassical transport loop}

\paragraph{What NEOPAX does.}
NEOPAX builds a reduced neoclassical transport model from a MONKES-style
database $D_{ij}(r,\nu,E_r)$ and evolves profiles with a JAX-native ODE solve
\citep{NEOPAXRepo}.

\paragraph{Governing equations.}
The thermal transport coefficients are assembled into an $L_{ij}$ matrix, from
which particle flux, heat flux, and parallel flow are computed:
\begin{equation}
\Gamma_a = -n_a\left(L_{11}A_1 + L_{12}A_2 + L_{13}A_3\right),
\end{equation}
\begin{equation}
Q_a = -T_a n_a\left(L_{21}A_1 + L_{22}A_2 + L_{23}A_3\right),
\end{equation}
\begin{equation}
U_{\parallel a}
=
-n_a\left(L_{31}A_1 + L_{32}A_2 + L_{33}A_3\right).
\end{equation}

The current implementation also forms an ambipolar electric-field source from
the charge-flux imbalance, schematically
\begin{equation}
S_{E_r} \propto -\Gamma_e + \Gamma_D + \Gamma_T,
\end{equation}
and integrates the resulting radial evolution equations with \texttt{diffrax}.

\paragraph{Inputs and outputs.}
Inputs are a MONKES-style $D_{ij}$ database, geometry, species profiles, and
edge conditions. Outputs include $E_r$, neoclassical particle and heat fluxes,
and profile evolution histories.

\paragraph{How NEOPAX fits in the workflow.}
It can replace a more expensive full transport loop when a differentiable,
reduced-order neoclassical model is preferable.

\section{Which quantities are optimization targets and which are transport inputs?}

This distinction is easy to miss and is crucial when building a black-box
workflow:

\begin{itemize}
  \item \textbf{Geometry metrics} such as aspect ratio, boundary smoothness, or target rotational transform are equilibrium-stage objectives.
  \item \textbf{Boozer-spectrum metrics} such as dominant symmetry-breaking $B_{mn}$ harmonics are usually optimization objectives and diagnostics.
  \item \textbf{Effective ripple} $\epsilon_{\mathrm{eff}}$ is usually a direct objective or screening metric, not a transport variable that Trinity3D advances.
  \item \textbf{SFINCS fluxes} are transport variables: they can directly feed Trinity3D or another transport solver.
  \item \textbf{GX / GENE heat fluxes} are both optimization objectives and transport inputs.
  \item \textbf{Bootstrap current} is a transport output that can become a new equilibrium input.
\end{itemize}

The user-supplied Mermaid chart reflects this correctly: some outputs are
promoted to the next physics stage, while others are used only as optimization
targets or screening diagnostics.

\section{Why the final temperature and density profiles differ from the initial guesses}

This is the central feedback mechanism in the workflow.

The equilibrium calculation requires a pressure and current profile, but those
profiles are often only a convenient starting point. Once neoclassical and
turbulent transport are evaluated, the radial particle and heat fluxes tell us
that the original profiles were not generally in balance with the sources and
losses.

The transport stage therefore updates:
\begin{itemize}
  \item species densities $n_s(r)$,
  \item species temperatures $T_s(r)$,
  \item the ambipolar electric field $E_r(r)$,
  \item and sometimes the bootstrap-current contribution.
\end{itemize}

The total pressure is then
\begin{equation}
p(r) = \sum_s n_s(r)T_s(r),
\end{equation}
so when the transport solver changes $n_s$ and $T_s$, the equilibrium pressure
profile changes too. If the bootstrap current changes appreciably, the current
profile also changes, which can shift the equilibrium again.

Thus a consistent design loop is
\begin{equation}
(R_{mn},Z_{mn},p,j)
\rightarrow
\text{equilibrium}
\rightarrow
\text{Boozer geometry}
\rightarrow
\text{transport fluxes}
\rightarrow
(n,T,E_r)
\rightarrow
(p_{\mathrm{new}},j_{\mathrm{new}})
\rightarrow
\text{equilibrium again}.
\end{equation}

\section{Examples from the literature}

The workflow described here is not hypothetical. Several published studies
already implement large parts of it:

\begin{itemize}
  \item \textbf{Precise quasi-symmetry optimization.}
  Landreman and Paul optimized stellarator magnetic fields for extremely
  precise quasi-symmetry, providing an archetypal example of equilibrium plus
  Boozer plus orbit-confinement analysis \citep{LandremanPaul2022}.

  \item \textbf{Direct turbulence optimization.}
  Kim et al.\ coupled DESC and GX directly inside an optimization loop, using
  nonlinear heat flux as an objective and Trinity3D for macroscopic profile
  evaluation \citep{Kim2024TurbulenceOpt}. This is very close to the workflow
  summarized in this repository.

  \item \textbf{First-principles profile prediction.}
  Bañón Navarro et al.\ coupled GENE-3D, KNOSOS, and a 1D transport solver to
  produce first-principles profile predictions for optimized stellarators
  \citep{BanonNavarro2023}. Although the specific neoclassical and transport
  codes differ, the logic is the same as SFINCS + GENE / GX + Trinity3D or
  NEOPAX.

  \item \textbf{Fast neoclassical loops for optimization.}
  MONKES was explicitly developed to make monoenergetic neoclassical
  calculations fast enough for optimization loops \citep{Escoto2024MONKES}.
\end{itemize}

Taken together, these studies show that the remaining challenge is primarily
software engineering: standard interfaces, reproducible containers, and an
orchestration layer.

\section{What a future containerized black-box workflow should standardize}

If the end goal is a one-click workflow, the most important design choice is
not the graphical interface but the contract between stages. The minimum useful
contracts are:

\begin{itemize}
  \item an \textbf{equilibrium container} that accepts boundary and profile coefficients and emits a VMEC-compatible \texttt{wout} plus a small machine-readable summary,
  \item a \textbf{Boozer container} that accepts \texttt{wout} and emits \texttt{boozmn} plus a summary of symmetry-breaking harmonics,
  \item a \textbf{neoclassical container} that accepts Boozer geometry and profile gradients and emits $\epsilon_{\mathrm{eff}}$, bootstrap current, $E_r$, particle flux, and heat flux in a shared HDF5 schema,
  \item a \textbf{turbulence container} that accepts geometry plus local profiles and emits heat and particle fluxes in a shared NetCDF or HDF5 schema,
  \item a \textbf{transport container} that accepts those fluxes plus source models and emits updated profiles and fusion metrics.
\end{itemize}

That, rather than the physics itself, is the key engineering task for turning a
research workflow into a black-box system.

\bibliographystyle{plainnat}
\bibliography{references}

\end{document}
